\documentclass[aspectratio=169,10pt,english]{beamer}

\input{../../../_preambulo_beamer.tex}
\input{../../../_comandos_beamer.tex}

\selectlanguage{english}

\title{MA1001B - Statistical Modeling for Decision Making}
\subtitle{Lesson 47: Simple Linear Regression: Least Squares and R-squared}
\author{}
\institute{Tecnol\'ogico de Monterrey}
\date{}

\begin{document}

\begin{frame}[plain]
  \vspace{1.2cm}
  {\Large\bfseries MA1001B - Statistical Modeling for Decision Making\par}
  \vspace{0.7cm}
  {\large Lesson 47: Simple Linear Regression: Least Squares and $R^2$\par}
  \vspace{0.7cm}
  {\color{TecRojo}\rule{\textwidth}{1.2pt}}
  \vspace{0.8cm}

  MA1001B Faculty\\[0.3cm]
  Term\\[0.3cm]
  Tecnol\'ogico de Monterrey
\end{frame}

\begin{frame}{The Line-Fitting Question}
  \begin{alertblock}{Guiding question}
    How can one numerical predictor be summarized by a unique line, and why
    is a high $R^2$ still not a causal coaching lever?
  \end{alertblock}

  \vspace{0.6em}
  By the end of this lesson, you can:
  \begin{itemize}
    \item read a scatter of coaching hours and units per hour;
    \item compute the least-squares slope and intercept;
    \item form $R^2=1-SSE/SST$;
    \item refuse to treat a high $R^2$ as a policy button.
  \end{itemize}

  \vfill
  \scriptsize All weeks and productivity values used today are synthetic. No real company data are used.
\end{frame}

\begin{frame}{Forty Synthetic Weeks}
  \small
  \begin{center}
    \begin{tabular}{lr}
      \toprule
      \textbf{Quantity} & \textbf{Value} \\
      \midrule
      Weeks $n$ & 40 \\
      Mean coaching hours $\bar x$ & 25.075277 \\
      Mean units per hour $\bar y$ & 23.315500 \\
      Sample correlation $r$ & 0.910813 \\
      True line & $y=12+0.45x+\varepsilon$ \\
      \bottomrule
    \end{tabular}
  \end{center}

  \vspace{0.5em}
  \begin{exampleblock}{Observational sample}
    $x$ is coaching hours that week. $y$ is units packed per labor hour.
    Coaching was observed, not randomly assigned.
  \end{exampleblock}
\end{frame}

\begin{frame}[fragile]{Step 1: Scatter, No Line Yet}
  \begin{columns}[T]
    \begin{column}{0.40\textwidth}
      \small
      \begin{lstlisting}
x = coaching_hours
y = units_per_hour
r = corrcoef(x, y)
      \end{lstlisting}

      \begin{block}{Verified output}
        $n=40$, $r=0.910813$\\
        $\bar x=25.075277$\\
        $\bar y=23.315500$
      \end{block}
    \end{column}
    \begin{column}{0.60\textwidth}
      \centering
      \includegraphics[width=\textwidth]{../figures/simple_regression_01_scatter.png}
      \vspace{0.2em}
      \scriptsize Scatter of 40 weeks from Step 1
    \end{column}
  \end{columns}
\end{frame}

\begin{frame}{Least Squares Minimizes Vertical Gaps}
  \small
  \[
    b_1=\frac{S_{xy}}{S_{xx}}=0.405936,
    \qquad
    b_0=\bar y-b_1\bar x=13.136554.
  \]
  Fitted line: $\hat y=13.136554+0.405936\,x$. True slope $0.45$; true
  intercept $12$. $SSE=113.102672$. scipy \texttt{linregress} matches the
  manual slope and intercept.
\end{frame}

\begin{frame}[fragile]{Step 2: The Fitted Line}
  \begin{columns}[T]
    \begin{column}{0.40\textwidth}
      \small
      \begin{lstlisting}
b1 = Sxy / Sxx
b0 = ybar - b1 * xbar
yhat = b0 + b1 * x
      \end{lstlisting}

      \begin{block}{Verified output}
        $b_1=0.405936$\\
        $b_0=13.136554$\\
        $SSE=113.102672$
      \end{block}
    \end{column}
    \begin{column}{0.60\textwidth}
      \centering
      \includegraphics[width=\textwidth]{../figures/simple_regression_02_least_squares.png}
      \vspace{0.2em}
      \scriptsize Least-squares line from Step 2
    \end{column}
  \end{columns}
\end{frame}

\begin{frame}{Step 3: High $R^2$ Is Not a Causal Lever}
  \begin{columns}[T]
    \begin{column}{0.42\textwidth}
      \small
      $SST=663.669939$, $SSR=550.567267$, $R^2=0.829580$.

      \vspace{0.4em}
      Fitted +10 hour gain: $4.059355$ units per hour.

      \vspace{0.5em}
      \textbf{Limit:} high $R^2$ does not imply a causal lever. Coaching was
      not randomized. Lesson 48 adds $SE(b_1)$, a $t$ test, and an
      influential point.
    \end{column}
    \begin{column}{0.58\textwidth}
      \centering
      \includegraphics[width=\textwidth]{../figures/simple_regression_03_r_squared_limit.png}
      \vspace{0.2em}
      \scriptsize $R^2$ decomposition from Step 3
    \end{column}
  \end{columns}
\end{frame}

\begin{frame}{Concept Check}
  \begin{enumerate}
    \item What vertical quantity does least squares minimize?
    \item Why can $b_1=0.405936$ differ from the true slope $0.45$?
    \item Verify $SSR+SSE=SST$ from the printed sums of squares.
    \item Why is $R^2=0.829580$ not a reason to mandate 10 more coaching hours?
  \end{enumerate}

  \vspace{0.6em}
  \begin{block}{Connection to Lesson 48}
    The session can continue directly into slope inference, residual plots,
    and influential points.
  \end{block}

  \vfill
  \scriptsize Source: OpenStax, \textit{Introductory Business Statistics 2e},
  Sections 13.3--13.4. CC BY-NC-SA.
\end{frame}

\appendix



\end{document}
